\subsection{与实极化几何量子化的比较}\label{branerp}

本节与下一节推广第 \ref{ccb} 节对膜量子化与几何量子化的比较。本节考虑一般情形：$Y$ 有全纯极化 $\varPi$，它限制为 $M$ 的实极化 $\P$；第 \ref{branecp} 节考虑 $\varPi$ 限制为复极化的情形。

所得图景与第 \ref{cotangent} 节的讨论非常相似，那里假定 $M=T^*N$。但此处结构比先前假定的更一般。极化 $\varPi$ 在 $M$ 上确定投影 $\pi:M\to N$，其纤维是 $\P$ 的叶。拓扑上可以选取嵌入 $\iota:N\to M$ 满足 $\pi\circ\iota=1$。一旦如此，$M$ 可等同于 $T^*N$。但与第 \ref{cotangent} 节不同，这里没有自然嵌入，也不能假定辛形式 $\omega_J$ 就是标准的 $\sum_i\d p_i\d q^i$。只能说：由于 $\pi$ 的纤维是拉格朗日的，$\omega_J$ 没有 $\d p_i\d p_j$ 项；由于假定 $\omega_J|_M$ 处处非退化，可选取纤维坐标 $p_i$ 使 $\d p\d q$ 项为通常的 $\sum_i\d p_i\d q^i$；但对 $\d q^i\d q^j$ 项，除它们不依赖 $p$（因 $\d\omega_J=0$）且闭之外，一无所知。故 $\omega_J$ 的一般形式为 $\omega_J=\sum_i\d p_i\d q^i+\pi^*(\alpha)$，$\alpha$ 为 $N$ 上的闭 2-形式。

这正是几何量子化中用实极化量子化辛流形 $M$ 时所假定的一般辛结构。故可望在更大的一般性下比较膜量子化与几何量子化。第 \ref{ccb} 节的分析假定了 $\alpha=0$。

几何量子化中，量子 Hilbert 空间 $\H_{\gq}$ 定义如下。取点 $w\in N$。预量子线丛 $\frL\to M$ 限制到 $\pi^{-1}(w)\cong\R^n$ 上平坦。故可定义线丛 $\frL_N\to N$，其在 $w$ 处的纤维是 $\frL$ 在 $\pi^{-1}(w)$ 上的协变常截面空间。此线丛没有自然联络，但有从 $\frL$ 的 Hermite 度量继承的自然 Hermite 度量。Hilbert 空间 $\H_{\gq}$ 定义为 $N$ 上取值于 $\frL_N$ 的 $\ltwo$ 半密度空间，内积为显然的
\be\label{nho}\langle\chi,\psi\rangle=\int_N\bar\chi\psi. \ee

为从膜量子化恢复这样的描述，我们完全仿照第 \ref{cotangent} 节，在 $\H'=\Hom(\B_\cc,\cF)$ 中定义一类候选 delta 函数态。如前，令 $F$ 为含给定点 $w\in N$ 的 $\varPi$ 叶，$\cF$ 为支撑在 $F$ 上、$\CP$ 丛平凡的 $(B,A,A)$ 型膜。并取与 $\alpha=(\d\vec p)^{1/2}$ 相应的 $\Hom(\B_\cc,\cF)$ 中元素。还需 $\Hom(\cF,\B)$ 中的一个角。如前，这是 $\L=\frL^{-1}$ 在 $F_M=F\cap M=\pi^{-1}(w)$ 上的协变常截面空间。令 $\beta$ 为这样的截面。它自然取值于 $\frL_N$ 在 $w$ 处纤维的对偶，故 $\beta^{-1}$ 取值于 $\frL_N$ 在 $w$ 处的纤维。

有了这些数据，图 \ref{strip} 的构造给出向量 $\psi'(w;\alpha,\beta)\in\H'$。由与先前相同的论证，两个这样的态在 $w\ne w'$ 时内积为零，这提示应把它们看作 delta 函数态。令 $\Psi$ 为 $\H$ 中任一态。内积 $f(w;\alpha,\beta)=(\psi'(w;\alpha,\beta),\Psi)$ 自然定义，且如式 (\ref{oko})，可定义 $N$ 上取值于 $\frL_N$ 的半密度
\be\label{rno} \h\Psi(w)=f(w;\alpha,\beta)\sqrt{\d q}\,\beta^{-1}. \ee
这给出从 $\Psi\in\H=\Hom(\B,\B_\cc)$ 到 $\h\Psi\in\H_{\gq}$ 的自然映射。我们预期它是同构。

作为用这种方式量子化的扭余切丛例子，可取 $M$ 为实代数流形 $x^2-y^2-z^2=j^2$，通过复化为由同一方程（复 $x,y,z$）定义的复流形 $Y'$（式 (\ref{nuuxx})）来量子化。$M$ 有两个实极化 $\P$ 与 $\P'$，由式 (\ref{horiz}) 中 $q,q'$ 取实定义。这两个极化由 $Y'$ 的两个全纯极化 $\varPi$ 与 $\varPi'$ 诱导，后者由同一公式（复变量）定义。故用 $\P$ 或 $\P'$ 对 $M$ 作几何量子化将给出等价的 Hilbert 空间。

为显式写出 $M$ 的这两个量子化，须描述 $N\cong\RP^1$ 上的线丛 $\frL_N$。由于 $M\cong T^*\RP^1\cong T^*S^1$ 非单连通，预量子线丛 $\frL\to M$（定义为带曲率 $\omega$ 的联络的酉线丛）不唯一；可通过把绕 $S^1$ 的和乐乘以任意 $\UU(1)$ 元素来修改。下面的选择是唯一使 $\PGL(2,\R)$ 在 $M$ 上的作用提升为同一群（而非其覆盖群）在 $\frL$ 上的作用的选择。

用两个开集 $\U$ 与 $\t\U$ 覆盖 $N$，$\U$ 由实变量 $q$ 参数化，$\t\U$ 由 $\t q=-1/q$ 参数化。辛形式为 $\omega=\d p\d q=\d\t p\d\t q$。$p$ 与 $\t p$ 的关系已在式 (\ref{anform}) 中得到；取 $\j=\i s$，该关系为
\be\label{relp} p=\t q^2\t p+\frac{2s}{q}. \ee
如下描述预量子线丛。在区域 $\U$ 中，$\frL$ 有平凡化使联络形式为 1-形式 $A=p\d q$；在区域 $\t\U$ 中，有平凡化给出 1-形式 $\t A=\t p\d\t q$。它们满足 $\d A=\d\t A=\omega$，符合 $\frL$ 上联络的预期。相应的协变微分为 $\d_A=\d+\i A$、$\d_{\t A}=\d+\i\t A$。两个联络形式相差一个规范变换：
\be\label{elp} A=\t A+2s\frac{\d q}{q}=\t A+2s\,\d\log q, \ee
故
\be \d_A=q^{-2\i s}\d_{\t A}q^{2\i s}. \ee
\begin{kanwu}[共轭方向]原文此处写的是 $\d_A=q^{2\i s}\d_{\t A}q^{-2\i s}$。按原文自己的约定 $\d_A=\d+\i A$ 与式 (\ref{elp})，直接计算给出上式；原文的共轭方向也与后文粘合律 $\psi(q)=q^{-2\i s}\t\psi(\t q)$ 不符。（若原文实际采用 $\d_A=\d-\i A$，则原文公式自洽。）此处按 $\d_A=\d+\i A$ 写出。\end{kanwu}

考虑 $\frL$ 的截面：在区域 $\U$ 中，相对于写 $A=p\d q$ 所用的平凡化，由函数 $\psi(q)$ 表示；在区域 $\t\U$ 中，相对于写 $\t A=\t p\d\t q$ 所用的平凡化，由函数 $\t\psi(\t q)$ 表示。式 (\ref{elp}) 提示在 $\U\cap\t\U$ 中它们应由 $\psi(q)=q^{-2\i s}\t\psi(\t q)$ 相联系。但须小心：$\U\cap\t\U$ 不连通，由两个分支组成，一个 $q>0$，一个 $q<0$。可把粘合律推广为 $\psi(q)=e^{\i\alpha}q^{-2\i s}\t\psi(\t q)$，$\alpha$ 为实常数。若对正负 $q$ 用同一常数，则无影响；若用不同常数，则有影响；这种可能性对应于一族单参数选择的 $\frL$。我们这样固定 $\frL$：对正负 $q$ 都要求 $\psi(q)=|q|^{-2\i s}\t\psi(\t q)$。

为方便描述，先注意测度 $\d q$ 满足
\be\label{measures}\frac{\d q}{q^2}=\d\t q.\ee
由于测度（与 1-形式不同）为正，可把此式提升到任意复幂。于是 $\psi$ 与 $\t\psi$ 的粘合律等价于
\be\label{convenone} \psi(q)(\d q)^{\i s}=\t\psi(\t q)(\d\t q)^{\i s}. \ee
故 $\frL\to N\cong\RP^1$ 的截面可理解为 $\RP^1$ 上权为 $\i s$ 的密度。

Hilbert 空间 $\H=\Hom(\B,\B_\cc)$ 由 $\RP^1$ 上取值于 $\frL$ 的半密度组成，故它是 $\RP^1$ 上 $\frac12+\i s$ 密度的空间。扭余切丛量子化中的 Hilbert 空间内积已由式 (\ref{nho}) 给出。如膜量子化中一贯的那样，代数 $\A=\Hom(\B_\cc,\B_\cc)$ 自然作用在 $\H$ 上。本例中此代数由式 (\ref{agrops}) 的算符 $t_-,t_+,t_3$ 生成。它们作为 $\H$ 上的算符都是反 Hermite 的。在 $\H$ 上，这些算符生成 $\PSL(2,\R)$ 的表示，对任意复 $s$ 都有定义。对实 $s$，此表示在 Hilbert 空间结构 (\ref{nho}) 下显然是酉的。

我们本可以从 $Y$ 的另一个全纯极化 $\varPi'$ 出发。那将导致同样的公式，只须作代换 $s\to-s$。故膜量子化预言：与权 $\frac12+\i s$ 和 $\frac12-\i s$ 的密度相应的 $\PSL(2,\R)$ 酉表示同构。这个同构是众所周知的，通常用显式初等公式证明，或仅注意两个表示都没有最高或最低权向量，且二次 Casimir 算符取值相同。

对更一般的 $\frL$ 选择，我们将得到 $\RP^1$ 上取值于某平坦线丛的 $\frac12+\i s$ 密度，该线丛绕 $\RP^1\cong S^1$ 一周有单值性 $e^{\i\veps}$，$\veps\in\R$。此单值性在拓扑上无关紧要；各线丛 $\frL_N\to N$ 作为酉线丛同构。但量子理论确实非平凡地依赖 $\veps$。同一代数 $\A$ 仍由同样的公式 (\ref{agrops}) 作用，与 $\veps$ 无关。但这些算符生成的表示依赖 $\veps$。对 $e^{\i\veps}=-1$，实 $s$ 时得到一族 $\SL(2,\R)$ 的酉表示；对其他 $\veps$，得到 $\PSL(2,\R)$ 的另一覆盖群（一般为通用覆盖）的一族表示。所有情形下等价关系 $s\leftrightarrow-s$ 仍成立。
